\documentclass[12pt]{book}
\usepackage{amsmath}
\usepackage{amssymb}
\usepackage{geometry}
\geometry{margin=1in}

\title{Volume 02: Field Theory, Mechanics, and Mathematical Structure \\ Book 02: Mathematical Operators and State Space}
\author{James C. Harvey}
\date{December 2025}

\begin{document}
\maketitle
\tableofcontents

\chapter{The 28-Dimensional State Space}

\section*{Abstract}
This book formalizes the SDT state space as a 28-dimensional operator manifold. It provides a narrative
derivation of each level of the state vector, expresses the occlusion operator as a geometric functional,
and demonstrates how force hierarchy emerges from the same state variables through regime selection.
All formulas are embedded in-line and mapped to the codebase definitions.

\section{From Space to State}
The state space is not an abstract convenience; it is a direct encoding of SDT’s physical grammar:
space, matter, movement, now. The matrix provides the substrate, flux provides the directional flow,
the medium defines collective behavior, and spations define the discrete lattice units. The 28-dimensional
state vector $\Xi \in \mathbb{R}^{28}$ is the minimal representation that captures geometry, dynamics,
and energetic manifestation.

\section{Hierarchical Structure of the 28 Components}
The state vector is built by levels:
\begin{itemize}
\item Level 1 (1): Existence
\item Level 2 (2): Line — position and relocation
\item Level 3 (3): Plane — boundary existence, planar relocation, rotation
\item Level 4 (4): Sphere — shell existence, shell relocation, rotation, orientation
\item Level 5 (5): Torus — matter structure topology
\item Level 6 (6): Dynamism — movement and ordering parameters
\item Level 7 (7): Energy — force manifestation modes
\end{itemize}

\section{State Vector Definition}
The 28 components are:
\begin{align}
\Xi =&\ (\xi_0; \xi_{10}, \xi_{11}; \xi_{p0}, \xi_{p1}, \xi_{p2}; \xi_{s0}, \xi_{s1}, \xi_{s2}, \xi_{s3}; \\
&\ T_1, T_2, T_3, T_4, T_5; \Phi_0, \Phi_1, \Phi_2, \Phi_3, \Phi_4, \Phi_5; \epsilon_0, \epsilon_1, \epsilon_2, \epsilon_3, \epsilon_b, \epsilon_4, \epsilon_5).
\end{align}
These values encode boundary scale, topology, movement, ordering, and energy configuration.

\section{Level-by-Level Meaning}
Each level is a closure of geometric intent:
\begin{itemize}
\item \textbf{Level 1} establishes existence as a boundary condition.
\item \textbf{Level 2} establishes location and relocation as flux-driven movement.
\item \textbf{Level 3} introduces planar boundary structure and rotational potential.
\item \textbf{Level 4} introduces shell structure and orientation of boundary massing.
\item \textbf{Level 5} encodes toroidal matter geometry.
\item \textbf{Level 6} encodes ordering, oscillation, and dynamical change.
\item \textbf{Level 7} encodes energetic manifestation as observable force modes.
\end{itemize}

\section{State as Narrative}
The state vector is the narrative of SDT written in numbers. It is not merely a coordinate system; it is
the algebra of space, matter, movement, and now.

\chapter{The Toroidal Layer: Matter Structure}

\section*{Abstract}
The toroidal layer ($T_1$–$T_5$) is the physical encoding of matter boundary geometry. This chapter
derives the role of each component and shows how they enter occlusion and coupling expressions.

\section{Component Definitions}
\begin{align}
T_1 &: \text{central ring scale},\\
T_2 &: \text{tube scale},\\
T_3 &: \text{topological surface area},\\
T_4 &: \text{polarized volume},\\
T_5 &: \text{aspect gradation (pressure slope)}.
\end{align}

\section{Boundary Scaling}
From $T_3$ we define the effective radius:
\begin{equation}
r_{\text{eff}} = \sqrt{\frac{T_3}{4\pi}},
\end{equation}
which governs solid-angle occlusion and therefore coupling strength.

\chapter{The Dynamical Layer: Ordering and Change}

\section*{Abstract}
The $\Phi$ components encode ordering and change. This layer is the mathematical representation of
movement and now.

\section{Ordering Parameters}
\begin{align}
\Phi_0 &: \text{omnidirectionality},\\
\Phi_1 &: \text{dynamic translocation},\\
\Phi_2 &: \text{oscillation frequency},\\
\Phi_3 &: \text{chirality indicator},\\
\Phi_4 &: \text{state variance},\\
\Phi_5 &: \text{transition potential}.
\end{align}
These parameters are the numerical grammar of change.

\chapter{State Operators and Conservation}

\section*{Abstract}
Operators define how the state vector evolves under flux. Conservation laws are expressed as invariances
of operator structure.

\section{Operator Form}
Let $\mathcal{L}$ be the state operator. Evolution is
\begin{equation}
\frac{d\Xi}{d\mathcal{N}} = \mathcal{L}(\Xi),
\end{equation}
where $\mathcal{N}$ is the ordering parameter of now.

\section{Linearized Evolution}
For small perturbations $\delta\Xi$ about a state $\Xi_0$,
\begin{equation}
\frac{d(\delta\Xi)}{d\mathcal{N}} = \mathbf{J}_{\mathcal{L}}(\Xi_0)\, \delta\Xi,
\end{equation}
where $\mathbf{J}_{\mathcal{L}}$ is the Jacobian of $\mathcal{L}$. Stability is determined by the
eigenvalues of $\mathbf{J}_{\mathcal{L}}$.

\section{Conservation}
Conservation appears when $\mathcal{L}$ leaves specific invariants unchanged. These invariants correspond
to stable boundary configurations and preserved flux pathways.

\chapter{Mapping to Implementation}

\section*{Abstract}
The state vector is not purely theoretical. It is implemented in code for investigation and validation.

\section{Implementation Mapping}
The 28 components are represented explicitly in code, ensuring that each geometric term is accessible for
numerical evaluation.

\chapter{Regime Selection as State Filtering}

\section*{Abstract}
Regime selection is state filtering. The same state vector yields different observed forces depending on
occlusion and curvature values.

\section{Filtering Rule}
Define regime selector:
\begin{equation}
\mathcal{R}(\Xi) = \text{argmax}\{\text{coupling weight}\},
\end{equation}
which chooses the effective regime by occlusion weighting.

\chapter{Toroidal Geometry Operators}

\section*{Abstract}
The torus components ($T_1$–$T_5$) encode boundary geometry, flux interception, and aspect gradation.
These quantities are directly used to compute occlusion, coupling strength, and regime classification.

\section{Effective Radius from Surface}
From the topological surface $T_3$, define the effective radius:
\begin{equation}
r_{\text{eff}} = \sqrt{\frac{T_3}{4\pi}}.
\end{equation}

\section{Occlusion Operator}
For two state vectors separated by distance $d$, the occlusion operator is:
\begin{equation}
E = \frac{\Omega_1 + \Omega_2}{4\pi}\left(1 + \tanh\left(\frac{|T_5|}{10^{10}}\right)\right),
\end{equation}
with
\begin{equation}
\Omega_i = \frac{r_{\text{eff},i}^2}{d^2}.
\end{equation}
This $E \in [0,1]$ directly determines coupling regime.

\section{Occlusion Scaling}
For $|T_5| \ll 10^{10}$, the hyperbolic tangent admits the linear approximation
\begin{equation}
\tanh\left(\frac{|T_5|}{10^{10}}\right) \approx \frac{|T_5|}{10^{10}},
\end{equation}
so the occlusion operator reduces to
\begin{equation}
E \approx \frac{\Omega_1 + \Omega_2}{4\pi}\left(1 + \frac{|T_5|}{10^{10}}\right).
\end{equation}
This is the low-slope geometric regime used in weakly graded boundaries.

\chapter{Force Hierarchy from State Operators}

\section*{Abstract}
The force hierarchy is a state-space consequence. By selecting different occlusion and curvature regimes,
the same interaction equation yields electromagnetic or gravitational limits.

\section{Coulomb and Gravity Regimes}
Let $E_C \approx 0$ and $E_G$ be the screened value. The effective coupling ratio is:
\begin{equation}
\frac{F_C}{F_G} = \frac{1 - E_C}{(1 - E_G)\kappa_{\text{screen}}}\cdot A_{\text{CMB}},
\end{equation}
where $\kappa_{\text{screen}} \approx 10^{-9}$ and $A_{\text{CMB}} \approx 10^{30}$ encode geometric screening
and flux amplification. This yields the observed hierarchy without introducing separate forces.

\chapter{Phase Space and Ordering Parameters}

\section*{Abstract}
The $\Phi$ components encode ordering and accessible state volume. This chapter defines accessible phase
space volume in terms of toroidal geometry and external influences.

\section{Accessible Phase Space}
Define the accessible phase space volume:
\begin{equation}
\mathcal{V}_{\Phi} = \ln\left(T_1 T_2^2\right) + \ln(1+|\Phi_4|) + \ln\left(1+\frac{|\Phi_5|}{10^{-20}}\right) + \ln(N_{\epsilon}),
\end{equation}
where $N_{\epsilon}$ counts active energy modes. Larger $\Phi_4$ implies greater external influence and
broader configuration accessibility.

\section{Differential Sensitivity}
The sensitivity to ordering parameters follows by differentiation:
\begin{equation}
\frac{\partial \mathcal{V}_{\Phi}}{\partial \Phi_4} = \frac{1}{1+|\Phi_4|},\quad
\frac{\partial \mathcal{V}_{\Phi}}{\partial \Phi_5} = \frac{1}{10^{-20} + |\Phi_5|}.
\end{equation}
These relations quantify how external influence and transition potential expand accessible state volume.

\chapter{Energy Mode Encoding}

\section*{Abstract}
The energy components $\epsilon_0$–$\epsilon_5$ encode potential, kinetic, rotational, field, binding,
flux, and transmission. These are not independent substances; they are structured modes of matrix flux.

\section{Mode Definitions}
\begin{align}
\epsilon_0 &: \text{Potential}, &
\epsilon_1 &: \text{Kinetic}, &
\epsilon_2 &: \text{Rotational},\\
\epsilon_3 &: \text{Field}, &
\epsilon_b &: \text{Binding}, &
\epsilon_4 &: \text{Flux}, &
\epsilon_5 &: \text{Transmission}.
\end{align}
The state vector therefore encodes both geometry and energetic manifestation.

\chapter{Operator Use in Investigations}

\section*{Abstract}
The 28D operators are not theoretical ornaments; they are embedded in investigation workflows. This
chapter demonstrates their application to orbital speed tests and empirical investigation templates.

\section{Operator Application}
The investigation pipeline constructs a state vector, computes occlusion $E$, evaluates hierarchy, and
maps predictions onto observed regimes.

\chapter{Narrative Closure}

\section*{Abstract}
State space is the mathematical spine of SDT. It encodes geometry, movement, ordering, and energy in a
single object. The narrative of SDT can therefore be read as the unfolding of $\Xi$ under matrix flux.

\section{From State to World}
All physical structure is the matrix expressed through state. The equations are not separate laws but
projections of one state onto different regimes. That is why the theory is one book: the state vector is
the condensed form of the world.

\end{document}
